% BSF Research Plan -- annotated skeleton.
%
% Every [[...]] is a slot to fill; grant_check.py reports each one left as a BLOCK, so
%   python3 <kit>/skills/grant-review/tools/grant_check.py --profile bsf research_plan.tex
% tells you what is still unwritten. Comments starting "% BSF §" point to the rule in
% skills/bsf-review/references/BSF_writing_guide.md. Delete the guidance comments when done
% if you like; they are never typeset.
%
% Format (BSF §3): 12 pt, line spacing 1.5 or 2, at most 38 lines per page, margins >= 2 cm,
% A4 (or US Letter if the initiating PI is in the U.S.). \linespread{1.5} gives ~34 lines;
% setspace's \onehalfspacing gives ~40 and breaks the cap. Items 1-5 fit in 15 pages,
% figures included; references go after them and do not count.
\documentclass[12pt,a4paper]{article}
\usepackage[margin=2cm]{geometry}
\linespread{1.5}
\usepackage[T1]{fontenc}
\usepackage{graphicx}
\usepackage{booktabs}
\usepackage{tabularx}
\usepackage[hidelinks]{hyperref}
\usepackage[font=small,labelfont=bf]{caption}   % captions may be smaller, but readable (BSF §3)
\setlength{\parskip}{0.3em}
\setlength{\parindent}{0pt}
% Optional: margin comments keyed to rules while you revise (\grscore{BSF 4.6}{...}).
% \usepackage{grantreview}

\begin{document}

% The heading must contain "Research Plan". A title is not needed (BSF §3).
\section*{Research Plan}

% ---------------------------------------------------------------------------------------
% OPENING -- 1/2 page, ~200 words (BSF §4.1). The adviser's summary of you starts here.
% Eight moves, one or two sentences each. Item 2 begins around page 5 in the mandated
% order, so a referee must have the question, hypothesis, objectives and the reason for
% two groups by the end of this box.
% ---------------------------------------------------------------------------------------
\noindent\fbox{\parbox{\dimexpr\linewidth-2\fboxsep-2\fboxrule}{%
% 1 Stakes: the larger problem the answer matters to. Not a history of the field.
[[Stakes, one sentence.]]
% 2 The question and why it is open: what specifically is unknown + the barrier (craft §2.2-2.3).
[[The question, and the barrier that has kept it open.]]
% 3 Long-term goal (the vision), one sentence only at the BSF (craft §3).
[[Long-term goal of the two groups.]]
% 4 This grant's aim and the central hypothesis (or central question / conjecture, BSF §4.12).
Our central hypothesis is that [[testable claim]].
% 5 Why it is testable now: the enabling result, preferably your own, and whose it is.
[[Enabling result and which group produced it.]]\par
% 6 Objectives, a line each, bold, numbered as in item 2. Two or three; four at most.
\textbf{Objective 1:} [[Determine whether ...]]\\
\textbf{Objective 2:} [[Identify / quantify ...]]\\
\textbf{Objective 3:} [[... (the riskiest one last)]]\par
% 7 Why these two groups: what each brings that the other lacks, and that the question needs both.
[[The Israeli group brings ...; the U.S. group brings ...; the question needs both because ...]]
% 8 Payoff: what will be known at the end that is not known now.
[[Payoff, one sentence.]]
}}

% ---------------------------------------------------------------------------------------
% ITEM 1 -- 2 to 2.5 pages (BSF §4.2). An argument, not a survey. Two to four sub-areas,
% each tied to an objective, each in four beats: known -> limitation -> barrier -> how this
% project moves past it. Cover BOTH groups' literatures at the same depth. Explain local
% (Israeli) sites, cohorts or datasets in a clause: 4 in 5 referees are not Israeli. The BSF
% picks referees from your references, so cite the field fairly and fully.
% ---------------------------------------------------------------------------------------
\section{Background and state of the art: [[your words]]}

\subsection{[[Sub-area A, tied to Objective 1]]}
[[What is known. Its limitation. The barrier that holds the field there. How this project
moves past it.]]

\subsection{[[Sub-area B, tied to Objective 2]]}
[[Four beats, as above. Name the closest prior work and what is different here
(the reviewer form's ``transformative ... not repetitive'').]]

% End item 1 on the gap paragraph: the specific unknown, why it is unknown, what changed.
\paragraph{The gap.} [[The specific unknown, why it has stayed unknown, and what now makes it
answerable.]]

% ---------------------------------------------------------------------------------------
% 1b PRELIMINARY RESULTS -- 3/4 to 1 page (BSF §4.3). Only the one or two results that
% motivate the WHOLE project; evidence for each objective goes in its block below (2-3 pages
% of evidence in all). From BOTH groups; one joint result if you have it. Each figure: a
% caption that states the finding, the group that produced it, n and the comparison. Say what
% the data do NOT settle -- that is the reason for the grant. Start-up: argue feasibility here
% instead of showing data (BSF §12). Theory: special cases, partial results (craft §14).
% ---------------------------------------------------------------------------------------
\subsection{Preliminary results}

\paragraph{[[Finding 1, stated as a result]] ([[which group]]).}
[[What was done, what it shows, and what it does not settle.]]
% \begin{figure}[t]\centering
% \includegraphics[width=.8\linewidth]{fig1}
% \caption{\textbf{[[The finding, as a sentence.]]} [[Group, n, comparison.]]}
% \end{figure}

\paragraph{[[Finding 2]] ([[the other group]]).}
[[What was done, what it shows, and what it does not settle.]]

% ---------------------------------------------------------------------------------------
% ITEM 2 -- 1 page (BSF §4.4). Two or three objectives, four at most. Each: a bold knowledge
% claim (not "to understand / to study", not a deliverable), its working hypothesis, one
% line of route, who leads and what the other group contributes, and its link to the others.
% Do NOT split objectives by country. Interdependent across the groups, but not blocking in
% sequence: each yields a result if another underperforms. Order them so risk rises.
% ---------------------------------------------------------------------------------------
\section{Objectives and significance}

\textbf{Objective 1: [[Determine whether ...]]} Working hypothesis: [[...]]. [[Route in one
sentence.]] Led by the [[Israeli / U.S.]] group, with [[what the other group contributes]].
[[What it gives Objective 2.]]

\textbf{Objective 2: [[...]]} Working hypothesis: [[...]]. [[Route.]] Led by [[...]], with
[[...]]. [[Link.]]

\textbf{Objective 3: [[...]]} Working hypothesis: [[...]]. [[Route.]] Led by [[...]], with
[[...]]. [[If this is the risky one: the project stands without it because ...]]

% One sentence on how the objectives depend on each other across the groups without blocking.
[[Each objective feeds the other group's ... ; each yields a result even if ...]]

% Significance: forward-looking, not the background again. What changes in the field when
% the objectives are met, with a "because" (craft §8). Criterion B. One sentence of social or
% economic value at most; the Broader Impacts Statement carries the rest (BSF §8).
\paragraph{Significance.} [[Which belief is confirmed or overturned, which method becomes
possible, which questions open -- because ...]]

% ---------------------------------------------------------------------------------------
% ITEM 3 -- about 8 to 9 pages: overview, one block per objective, collaboration (BSF §4.5-4.7).
% At least as long as items 1 and 2 together. The regulation requires the respective roles
% of the Israeli and American PIs here.
% ---------------------------------------------------------------------------------------
\section{Methodology and plan of operation}

% Overview, 1/2 to 3/4 page (BSF §4.5): the logic of the plan, ONE design figure with the two
% groups as lanes and arrows for what crosses between them, shared methods stated once, and
% the timeline.
[[One paragraph: how the objectives follow from each other, and why this structure fits the
question.]]
% \begin{figure}[t]\centering
% \includegraphics[width=\linewidth]{design}
% \caption{\textbf{Project design.} Objectives as rows, the two groups as lanes; arrows are
% samples, data, models and people crossing between the groups.}
% \end{figure}

\paragraph{Shared methods.} [[Datasets, statistics, validation, data sharing between sites --
stated once here.]]

% Timeline: milestones, decision points and exchange visits marked; who works in each period.
% A milestone is a RESULT ("Month 18: rates for >= 30 communities; decides route A or B"),
% not an activity ("complete Objective 2 experiments"). It must tell the same story as the
% work schedule in the online form (criterion E).
\paragraph{Timeline.}
\begin{center}
\begin{tabularx}{\linewidth}{lXXXX}
\toprule
 & Year 1 & Year 2 & Year 3 & Year 4 \\
\midrule
Objective 1 & [[tasks, who]] & [[milestone]] & & \\
Objective 2 & & [[tasks, who]] & [[decision point]] & \\
Objective 3 & & & [[tasks, who]] & [[...]] \\
Visits & [[who, where, how long]] & [[...]] & [[...]] & [[...]] \\
\bottomrule
\end{tabularx}
\end{center}

% One block per objective, 2-2.5 pages each (BSF §4.6). The same five parts in every block.
% Spend method detail where the risk is: new or untested methods in full, routine ones briefly.
\subsection{Objective 1: [[the bold claim again, shortened]]}

% (1) Rationale and working hypothesis, 2-4 sentences.
\paragraph{Rationale and hypothesis.} [[Why this objective is needed and what it tests.]]

% (2) Preliminary evidence specific to this objective, and which group produced it.
\paragraph{Preliminary evidence.} [[A number or a figure for this objective, or a pointer to
the preliminary results in Section 1; which group produced it.]]

% (3) Tasks. Each: purpose, approach at expert level (controls, n or power, validation,
% analysis, data, how data move between the sites), the check that tells you it worked,
% and a lead line naming the group and what it receives from the other (BSF §4.6, §7).
\paragraph{Task 1.1: [[name]].} [[Purpose. Approach at the level a referee in this method
expects. The check.]] Led by the [[...]] group; [[what the other group supplies]].

\paragraph{Task 1.2: [[name]].} [[...]] Led by the [[...]] group; [[...]].

% (4) Expected outcomes and success criterion, stated before the work is done, and what each
% plausible outcome would mean. Theory: which statement or special case. Computation: which
% gain over which baseline on which held-out data, read as a finding.
\paragraph{Expected outcomes and success criterion.} [[The hypothesis is supported if ...;
it is rejected if ....]]

% (5) Pitfalls and alternatives, numbered if-then items, 10-20% of the block. Each: the
% likely problem, why it is unlikely given the evidence, the concrete alternative, why the
% objective survives. Both kinds: technical (the method fails) and interpretive (another
% explanation, and what separates them). End with what is learned even if the hypothesis fails.
\paragraph{Pitfalls and alternatives.} (i) If [[technical problem]] ([[why unlikely]]), then
[[alternative]], which still [[...]]. (ii) If the result could also be explained by
[[alternative explanation]], [[the control or comparison that separates them]]. Even if
[[the hypothesis fails]], we obtain [[...]].

\subsection{Objective 2: [[...]]}
\paragraph{Rationale and hypothesis.} [[...]]
\paragraph{Preliminary evidence.} [[...]]
\paragraph{Task 2.1: [[name]].} [[...]] Led by the [[...]] group; [[...]].
\paragraph{Expected outcomes and success criterion.} [[...]]
\paragraph{Pitfalls and alternatives.} (i) If [[...]], then [[...]].

% The last block may carry one bounded moonshot, marked as such (craft §9.4).
\subsection{Objective 3: [[...]]}
\paragraph{Rationale and hypothesis.} [[...]]
\paragraph{Preliminary evidence.} [[...]]
\paragraph{Task 3.1: [[name]].} [[...]] Led by the [[...]] group; [[...]].
\paragraph{Expected outcomes and success criterion.} [[...]]
\paragraph{Pitfalls and alternatives.} (i) If [[...]], then [[...]].

% Collaboration, 3/4 to 1 page (BSF §4.7, tests in §7). Criterion C, a separate gate.
\subsection{The collaboration and the roles of the Israeli and U.S. groups}

\paragraph{Why this pairing.} [[What only the Israeli group brings. What only the U.S. group
brings. Why the question needs both.]]

\begin{center}
\begin{tabularx}{\linewidth}{Xcc}
\toprule
Task & [[Israeli group]] & [[U.S. group]] \\
\midrule
1.1 [[name]] & lead & contributes \\
1.2 [[name]] & contributes & lead \\
2.1 [[name]] & co-lead & co-lead \\
\bottomrule
\end{tabularx}
\end{center}

\paragraph{Interface tasks.} [[The one or two co-led tasks where one group's output is the
other's input, with the month of each hand-off.]]

\paragraph{How the collaboration runs.} [[Exchange visits: who, how long (more than two full
working days), which year. Meeting cadence. Shared repository. Co-supervised students. The
joint papers you expect.]]

% If the PIs held a BSF grant together, say what it produced; one without joint papers must
% be addressed (BSF §7, §10). Governance in two or three sentences.
\paragraph{History and governance.} [[Previous joint work and what this grant adds. Who
decides on changes of direction; how authorship is settled.]]

% ---------------------------------------------------------------------------------------
% ITEM 4 -- 1/2 page (BSF §4.8). Only the risks that cross objectives; the per-objective
% fallbacks are in the blocks above -- point to them, do not repeat them.
% ---------------------------------------------------------------------------------------
\section{Risk analysis and alternative paths}

\begin{center}
\begin{tabularx}{\linewidth}{XllX}
\toprule
Risk & Likelihood & Effect & Mitigation \\
\midrule
{}[[Hand-off between the groups is late]] & [[...]] & [[...]] & [[each group works on archived or simulated inputs until it arrives]] \\
{}[[Shared instrument, permit, field season, recruitment]] & [[...]] & [[...]] & [[...]] \\
\bottomrule
\end{tabularx}
\end{center}

\paragraph{If the central hypothesis is wrong.} [[What the project still establishes, and why
the risk is worth taking.]]

% ---------------------------------------------------------------------------------------
% ITEM 5 -- 1/2 to 3/4 page (BSF §4.9). Shows the whole project is resourced; the grant mostly
% pays for the cooperation. Both sides. Consistent with the CV form's current grants and with
% the budget. Flag anything one side provides to the other (criterion C(b)).
% ---------------------------------------------------------------------------------------
\section{Available resources}

\paragraph{Israel ([[institution]]).} Personnel: [[name, role, effort, who pays]]. Equipment
and facilities: [[named instruments, compute, sites, datasets]]. Other funding: [[grant, and
why it does not overlap]].

\paragraph{U.S. ([[institution]]).} Personnel: [[...]]. Equipment and facilities: [[...]].
Other funding: [[...]].

\paragraph{Provided to the partner.} [[Facility, material or instrument one side makes
available to the other, and when.]]

% ---------------------------------------------------------------------------------------
% CLOSING -- 1/4 page (BSF §4.10). Not risk, not resources: what will be known, why these two
% groups now, what it opens next. Then the generative-AI note (BSF §3).
% ---------------------------------------------------------------------------------------
\section*{Outlook}
[[Three or four sentences.]]

\paragraph{Use of generative AI.} [[None was used. / Generative AI was used only for ...;
the PIs checked all content.]]

% References: here, after item 5, if you tick the checkbox; or as a separate upload headed
% "References". Full author lists up to four, else three + et al.; the BSF finds referees
% here (BSF §10).
% \bibliographystyle{unsrt}
% \bibliography{refs}

\end{document}
